\chapter{การพัฒนาและการประยุกต์ใช้ห้องปฏิบัติการเสมือนสะเต็มศึกษาขั้นสูง}
\label{chap:ch07}

\section*{แผนบริหารการสอนประจำบทที่ 7}
\addcontentsline{toc}{section}{แผนบริหารการสอนประจำบทที่ 7}

\noindent\textbf{หัวข้อเนื้อหาประจำบท}
\begin{enumerate}
    \item สถาปัตยกรรมการออกแบบห้องปฏิบัติการทัศนศาสตร์เสมือน (Virtual Optics Bench)
    \item การจำลองปรากฏการณ์คลื่นนิ่งและการแทรกสอด (Wave Interference Simulator)
    \item ห้องปฏิบัติการวงจรอิเล็กทรอนิกส์และออสซิลโลสโคปเชิงพื้นที่สามมิติ
    \item การสร้างระบบจำลองวงโคจรระบบดาวคู่ทางดาราศาสตร์ฟิสิกส์และการเผยแพร่ข้ามแพลตฟอร์ม
\end{enumerate}

\noindent\textbf{วัตถุประสงค์เชิงพฤติกรรม}
\begin{enumerate}
    \item พัฒนาระบบคำนวณการหักเหของแสงและการสร้างภาพด้วยเลนส์แบบเรียลไทม์ได้
    \item สร้างแบบจำลองการซ้อนทับของคลื่นสองมิติในถังคลื่นเสมือนจริงได้
    \item ออกแบบเครื่องมือวัดทางอิเล็กทรอนิกส์เชิงพื้นที่ที่เชื่อมโยงกับทฤษฎีวงจร RLC ได้
    \item บูรณาการและส่งมอบ (Deploy) แอปพลิเคชันห้องทดลองเสมือนสู่แพลตฟอร์มสากลได้
\end{enumerate}

\newpage

\noindent\textbf{กิจกรรมการเรียนการสอน}
\begin{enumerate}
    \item การสัมมนาเชิงปฏิบัติการวิเคราะห์สถาปัตยกรรมห้องทดลองวิทยาศาสตร์เสมือนจริง
    \item โครงงานกลุ่มพัฒนาห้องปฏิบัติการเสมือนจริงตามหัวข้อสะเต็มศึกษาที่กำหนด
    \item การทดสอบประสิทธิภาพการทำงานข้ามแพลตฟอร์ม (WebXR, PC VR, และ Standalone XR)
    \item การนำเสนอผลงานโครงงานและการประเมินความพึงพอใจของผู้ใช้งาน
\end{enumerate}

\noindent\textbf{สื่อการเรียนการสอน}
\begin{enumerate}
    \item เอกสารคู่มือการพัฒนาโครงงานสะเต็มศึกษาเสมือนจริงขั้นสูง
    \item ชุดโมเดลสามมิติอุปกรณ์วิทยาศาสตร์ความละเอียดสูง (PBR Scientific Assets)
    \item รหัสต้นฉบับโครงงานห้องปฏิบัติการทัศนศาสตร์ คลื่น และวงจรไฟฟ้า
    \item แพลตฟอร์มคลาวด์สำหรับโฮสต์และเผยแพร่ WebXR ผ่าน GitHub Pages CDN
\end{enumerate}

\noindent\textbf{การประเมินผล}
\begin{enumerate}
    \item การประเมินโครงงานพัฒนานวัตกรรมห้องปฏิบัติการเสมือนจริง (Term Project Rubrics)
    \item การทดสอบการทำงานของระบบข้ามแพลตฟอร์มโดยปราศจากข้อผิดพลาด
    \item การประเมินทักษะการสื่อสารและการนำเสนอผลงานทางวิชาการ
\end{enumerate}

\newpage

\section{ห้องปฏิบัติการทัศนศาสตร์เสมือนจริง}

การทดลองทัศนศาสตร์เป็นหนึ่งในตัวอย่างคลาสสิกที่ได้รับประโยชน์อย่างมหาศาลจากเทคโนโลยี XR เนื่องจากผู้เรียนสามารถมองเห็นลำแสงเลเซอร์ (Ray Tracing) และเส้นแนวแกนทัศน์ได้อย่างชัดเจนในมิติสามมิติ ซึ่งในการทดลองจริงไม่สามารถมองเห็นเส้นทางเดินแสงได้อย่างชัดเจนหากไม่มีควันหรือฝุ่นละออง

\subsection{สมการช่างทำเลนส์และการหักเหของแสง}

การคำนวณทางเดินแสงผ่านเลนส์บางสอดคล้องกับสมการช่างทำเลนส์ (Lensmaker's Equation)
\begin{equation}
\frac{1}{f} = (n - 1)\left(\frac{1}{R_1} - \frac{1}{R_2}\right)
\end{equation}
และความสัมพันธ์ระหว่างระยะวัตถุ $s$ ระยะภาพ $s'$ และความยาวโฟกัส $f$ เป็นไปตามสมการเกาส์เซียน
\begin{equation}
\frac{1}{s} + \frac{1}{s'} = \frac{1}{f}
\end{equation}
โดยมีกำลังขยายขวาง (Transverse Magnification)
\begin{equation}
m = -\frac{s'}{s} = \frac{y'}{y}
\end{equation}

\begin{figure}[htbp]
\centering
\begin{tikzpicture}[scale=1.0]
    % Optical Axis
    \draw[darkNavy, line width=1.0pt, dashed] (-4.5,0) -- (5.5,0) node[right]{\textbf{แกนมุขสำคัญ}};

    % Lens Plane
    \draw[spatialBlue, ultra thick] (0,-2.5) -- (0,2.5);
    \draw[spatialBlue, line width=1.5pt, fill=spatialBlue!25] (0,0) ellipse (0.3cm and 2.5cm);
    \node[spatialBlue, above] at (0, 2.5) {\textbf{เลนส์นูน}};

    % Focal Points
    \fill[amberGlow] (-2.5,0) circle (3pt) node[below]{\textbf{$F$}};
    \fill[amberGlow] (2.5,0) circle (3pt) node[below]{\textbf{$F'$}};

    % Object Arrow
    \draw[-{Stealth}, line width=2.0pt, openxrTeal] (-4.0,0) -- (-4.0,1.8) node[above]{\textbf{วัตถุ ($y$)}};
    \node[below] at (-4.0,0) {$-s$};

    % Principal Rays
    % Ray 1: Parallel to axis then through focal point
    \draw[red!90!black, line width=1.3pt] (-4.0,1.8) -- (0,1.8) -- (4.0,-1.2);
    % Ray 2: Through optical center
    \draw[emeraldActive, line width=1.3pt] (-4.0,1.8) -- (0,0) -- (4.0,-1.8);
    % Ray 3: Through front focus then parallel
    \draw[cyberPurple, line width=1.3pt] (-4.0,1.8) -- (-2.5,0) -- (0,-1.8) -- (4.0,-1.8);

    % Image Arrow
    \draw[-{Stealth}, line width=2.0pt, red!80!black] (3.33,0) -- (3.33,-1.5) node[below]{\textbf{ภาพจริง ($y'$)}};
    \node[above] at (3.33,0) {$s'$};
\end{tikzpicture}
\caption{การติดตามรังสีแสงสามมิติในระบบจำลองทัศนศาสตร์เสมือนจริง}
\label{fig:virtual_optics}
\end{figure}

\begin{workedexamplebox}{การคำนวณตำแหน่งและลักษณะภาพในห้องทดลองเสมือน}
วางวัตถุไว้ที่ระยะห่าง $s = 30.0\text{ cm}$ ด้านหน้าเลนส์นูนที่มีความยาวโฟกัส $f = 10.0\text{ cm}$ จงคำนวณระยะภาพ $s'$ กำลังขยาย $m$ และระบุชนิดของภาพที่จะต้องสร้างขึ้นในเอนจินกราฟิกส์

\textbf{วิธีทำ:}
จากสมการของเลนส์บาง
\begin{align}
\frac{1}{s'} &= \frac{1}{f} - \frac{1}{s} = \frac{1}{10.0} - \frac{1}{30.0} = \frac{3 - 1}{30.0} = \frac{2}{30.0} = \frac{1}{15.0}\text{ cm}^{-1} \\
s' &= +15.0\text{ cm}
\end{align}
คำนวณกำลังขยาย
\begin{equation}
m = -\frac{s'}{s} = -\frac{15.0}{30.0} = -0.5
\end{equation}
ผลลัพธ์แสดงว่าภาพที่เกิดขึ้นเป็น \textbf{ภาพจริง (Real Image)} อยู่หลังเลนส์ที่ระยะ $15.0\text{ cm}$ มีลักษณะ \textbf{หัวกลับ} ($m < 0$) และมีขนาดเป็น \textbf{ครึ่งหนึ่งของวัตถุจริง} ($|m| = 0.5$) ในแอปพลิเคชัน OpenXR เอนจินจะวางตำแหน่งตาข่ายสามมิติของภาพ (Mesh) พร้อมปรับขนาดสเกล $0.5$ เท่าและหมุนแกน $180^\circ$ โดยอัตโนมัติ
\end{workedexamplebox}

\section{การจำลองวงโคจรระบบดาวคู่ทางดาราศาสตร์ฟิสิกส์}

ในการศึกษาขั้นสูง การประยุกต์ OpenXR ร่วมกับแบบจำลองทางดาราศาสตร์ฟิสิกส์ เช่น การโคจรของระบบดาวคู่แบบแตะกัน (Contact Binary Stars) สอดคล้องกับงานวิจัยของ ทัศนา และคณะ ในการวิเคราะห์การถ่ายเทมวลและพลังงานผ่านผิวสมศักดิ์ของโรช (Roche Equipotential Surfaces)
\begin{equation}
\Phi(\mathbf{r}) = -\frac{G M_1}{r_1} - \frac{G M_2}{r_2} - \frac{1}{2}\omega^2 (x^2 + y^2)
\end{equation}
การจำลองแบบสามมิติช่วยให้นักศึกษาสามารถหมุนสังเกตจุดลากรองจ์ (Lagrange Points $L_1, L_2, L_3, L_4, L_5$) และการบิดเบี้ยวของชั้นบรรยากาศดาวในสภาพแวดล้อมเสมือนจริงได้อย่างลึกซึ้ง

\begin{table}[htbp]
\centering
\caption{เมทริกซ์การประยุกต์ใช้ห้องปฏิบัติการเสมือนจริงในสาขาวิชาต่าง ๆ}
\label{tab:stem_matrices}
\begin{tabularx}{\textwidth}{lYY}
\toprule
\textbf{สาขาวิชา} & \textbf{ปรากฏการณ์ที่จำลองใน XR} & \textbf{การปฏิสัมพันธ์ของผู้เรียนผ่าน OpenXR} \\
\midrule
ฟิสิกส์ทัศนศาสตร์ & ทางเดินแสง การหักเห การสะท้อนกลับหมด & ใช้มือเลื่อนตำแหน่งเลนส์ กระจก และฉากรับภาพ \\
ฟิสิกส์คลื่น & การแทรกสอดสองมิติ และคลื่นนิ่งในท่อ & ปรับความถี่และความยาวคลื่นด้วยปุ่มหมุนเสมือน \\
วิศวกรรมไฟฟ้า & วงจรเรโซแนนซ์ RLC และสัญญาณความถี่ & ใช้สายวัดต่อจุดวงจรและสังเกตกราฟออสซิลโลสโคป \\
ดาราศาสตร์ฟิสิกส์ & วงโคจรเคปเลอร์ และระบบดาวคู่แบบแตะกัน & ปรับอัตราส่วนมวลและมุมเอียงของการสังเกต \\
\bottomrule
\end{tabularx}
\end{table}

\section{ห้องปฏิบัติการวงจรอิเล็กทรอนิกส์และการหมุนเฟสเซอร์ 3 มิติ}

ในการศึกษาวิศวกรรมไฟฟ้าและฟิสิกส์คลื่นแม่เหล็กไฟฟ้า การทำความเข้าใจวงจรไฟฟ้ากระแสสลับ (Alternating Current: AC) เป็นเรื่องท้าทายสำหรับผู้เรียนเนื่องจากแรงดันไฟฟ้าและกระแสไฟฟ้ามีการเปลี่ยนแปลงเชิงมุมเฟสตลอดเวลา การประยุกต์ใช้การประมวลผลเชิงพื้นที่ 3 มิติช่วยให้นักศึกษาสามารถมองเห็น \textbf{แผนภาพเฟสเซอร์ (Phasor Diagram)} หมุนรอบจุดกำเนิดในปริภูมิเสมือนจริงได้อย่างเป็นรูปธรรม

\subsection{คณิตศาสตร์ของเฟสเซอร์และอิมพีแดนซ์ในวงจร RLC}

สัญญาณรูปคลื่นไซน์ของแรงดันไฟฟ้ากระแสสลับ $v(t) = V_m \cos(\omega t + \phi)$ สามารถแทนได้ด้วยเวกเตอร์หมุน (Phasor) ในระนาบเชิงซ้อน
\begin{equation}
\mathbf{V} = V_{\text{rms}} e^{j\phi} = V_{\text{rms}} (\cos\phi + j\sin\phi) \label{eq:phasor_def}
\end{equation}
โดยที่ $\omega = 2\pi f$ คือความถี่เชิงมุมของแหล่งกำเนิด

ในวงจรอนุกรม $RLC$ ที่ประกอบด้วยตัวต้านทาน ($R$), ตัวเหนี่ยวนำ ($L$), และตัวเก็บประจุ ($C$) ความสัมพันธ์ของเฟสเซอร์ระหว่างแรงดันตกคร่อมแต่ละอุปกรณ์แสดงได้ดังนี้:
\begin{enumerate}
    \item \textbf{ตัวต้านทาน ($R$):} แรงดัน $\mathbf{V}_R$ มีเฟสตรงกับกระแส $\mathbf{I}$ ($\phi = 0$)
    \item \textbf{ตัวเหนี่ยวนำ ($L$):} แรงดัน $\mathbf{V}_L$ มีเฟสนำหน้ากระแส $90^\circ$ ($+\frac{\pi}{2}\text{ rad}$) โดยมีค่ารีแอกแทนซ์ความเหนี่ยวนำ $X_L = \omega L$
    \item \textbf{ตัวเก็บประจุ ($C$):} แรงดัน $\mathbf{V}_C$ มีเฟสตามหลังกระแส $90^\circ$ ($-\frac{\pi}{2}\text{ rad}$) โดยมีค่ารีแอกแทนซ์ความจุ $X_C = \frac{1}{\omega C}$
\end{enumerate}

อิมพีแดนซ์รวมของวงจร (Total Impedance) $Z$ และมุมต่างเฟสรวม $\phi$ คำนวณได้จาก
\begin{align}
\mathbf{Z} &= R + j(X_L - X_C) \label{eq:impedance_complex} \\
Z &= |\mathbf{Z}| = \sqrt{R^2 + (X_L - X_C)^2} \label{eq:impedance_mag} \\
\phi &= \arctan\left(\frac{X_L - X_C}{R}\right) \label{eq:phase_angle}
\end{align}

\begin{figure}[htbp]
\centering
\begin{tikzpicture}[scale=1.1]
    % Complex Plane Axes
    \draw[-{Stealth}, line width=1.2pt, slateGrey] (-1.0,0) -- (5.0,0) node[right]{\textbf{แกนจริง (Real)}};
    \draw[-{Stealth}, line width=1.2pt, slateGrey] (0,-2.5) -- (0,3.5) node[above]{\textbf{แกนจินตภาพ (Imaginary)}};

    % Current Vector I along real axis
    \draw[-{Stealth}, line width=1.8pt, darkNavy] (0,0) -- (3.5,0) node[below]{\textbf{$\mathbf{I}$ / $\mathbf{V}_R$}};

    % Inductive Voltage VL (+j)
    \draw[-{Stealth}, line width=1.8pt, emeraldActive] (0,0) -- (0,2.8) node[left]{\textbf{$\mathbf{V}_L = j X_L \mathbf{I}$}};

    % Capacitive Voltage VC (-j)
    \draw[-{Stealth}, line width=1.8pt, amberGlow] (0,0) -- (0,-1.5) node[left]{\textbf{$\mathbf{V}_C = -j X_C \mathbf{I}$}};

    % Net Reactive Voltage (VL - VC)
    \coordinate (Vnet) at (0, 1.3);
    \draw[dashed, red!80!black, line width=1.0pt] (0, 1.3) -- (3.5, 1.3);
    \draw[dashed, slateGrey, line width=1.0pt] (3.5, 0) -- (3.5, 1.3);

    % Total Voltage V_total
    \draw[-{Stealth}, line width=2.2pt, spatialBlue] (0,0) -- (3.5,1.3) node[above right]{\textbf{$\mathbf{V}_{\text{total}}$}};

    % Phase Angle phi
    \draw[->, line width=1.2pt, cyberPurple] (1.2,0) arc (0:20.4:1.2) node[midway, right]{\textbf{$\phi$}};
    \node[red!80!black, left] at (0, 0.65) {$V_L - V_C$};
\end{tikzpicture}
\caption{แผนภาพเฟสเซอร์ในระนาบเชิงซ้อนจำลอง 3 มิติ แสดงความสัมพันธ์ระหว่างแรงดันและความถี่}
\label{fig:phasor_diagram}
\end{figure}

\section{การเผยแพร่และการส่งมอบสู่แพลตฟอร์มสากล}

เพื่อให้สื่อการเรียนรู้สะเต็มศึกษาที่สร้างขึ้นสามารถเข้าถึงได้อย่างกว้างขวาง สถาปัตยกรรม OpenXR รองรับการคอมไพล์และส่งมอบข้ามระบบ (Cross-Platform Deployment)
\begin{enumerate}
    \item \textbf{Standalone Headsets (Android XR / Meta Quest / Pico):} คอมไพล์เป็น APK ผ่าน Android NDK และ OpenXR Loader
    \item \textbf{PC VR (Windows / Linux Monado):} รันผ่าน Direct-to-Device Mode ด้วยประสิทธิภาพกราฟิกสูงสุด
    \item \textbf{WebXR via Global CDN:} รันบนเว็บเบราว์เซอร์ผ่านมาตรฐาน W3C WebXR Device API โดยไม่ต้องติดตั้งโปรแกรมเพิ่มเติม
\end{enumerate}

\begin{aipromptbox}[การสร้างห้องทดลองทัศนศาสตร์เสมือนจริง WebXR แบบ Standalone ด้วย AI]
\textbf{วัตถุประสงค์:} สร้างไฟล์ HTML/JavaScript แบบ Standalone ชิ้นเดียว (Single-file WebXR Application) สำหรับจำลองการติดตามรังสีแสงผ่านเลนส์นูน (Ray Tracing) แบบ 3 มิติ ที่สามารถเปิดรันบนเบราว์เซอร์ของแว่น XR ได้ทันทีโดยไม่ต้องพึ่งพาเซิร์ฟเวอร์ภายนอก\\
\textbf{โครงสร้าง CLEAR Prompt:}
\begin{itemize}[topsep=1pt, itemsep=1pt, leftmargin=1.2em]
    \item \textbf{Context (บริบท):} ระบบห้องทดลองฟิสิกส์เสมือนจริงสำหรับโรงเรียนมัธยมศึกษาและมหาวิทยาลัยที่รองรับมาตรฐาน W3C WebXR
    \item \textbf{Logic (ตรรกะ):} ใช้ Three.js ผ่าน CDN สร้างรางทัศนศาสตร์ เลนส์นูนบาง และฉากรับภาพ คำนวณรังสี 3 เส้นหลัก (ขนานแกน, ผ่านจุดกึ่งกลางเลนส์, ผ่านจุดโฟกัส) ด้วยสมการเกาส์เซียน $\frac{1}{s} + \frac{1}{s'} = \frac{1}{f}$
    \item \textbf{Expectation (ผลลัพธ์ที่ต้องการ):} โค้ด HTML สมบูรณ์ที่มีปุ่ม "Enter VR", ตัวเลื่อน (Sliders) ปรับระยะวัตถุ $s$ และความยาวโฟกัส $f$ พร้อมแสดงค่าระยะภาพ $s'$ และกำลังขยาย $m$ บนจอ HUD
    \item \textbf{Action \& Restriction (ข้อกำหนด):} ห้ามใช้ npm หรือ build tool อื่น ต้องรันได้ทันทีเมื่อบันทึกเป็น \texttt{index.html} และรักษาอัตราเฟรมคงที่ที่ 60 FPS
\end{itemize}
\end{aipromptbox}

\begin{codebox}[ระบบติดตามรังสีแสงผ่านเลนส์นูน 3 มิติในภาษา JavaScript (Three.js)]
// ฟังก์ชันคำนวณและอัปเดตเส้นทางเดินแสงในห้องทดลองเสมือน
function updateOpticsBench(objectDist, focalLength, objectHeight) {
    // คำนวณระยะภาพตามสมการ 1/s' = 1/f - 1/s
    const invImageDist = (1.0 / focalLength) - (1.0 / objectDist);
    const imageDist = 1.0 / invImageDist;
    const magnification = -imageDist / objectDist;
    const imageHeight = objectHeight * magnification;

    // จุดพิกัดสำคัญ (x=0 คือตำแหน่งเลนส์, แนวแกน Z คือแกนมุขสำคัญ)
    const pObjTip = new THREE.Vector3(0, objectHeight, -objectDist);
    const pLensCenter = new THREE.Vector3(0, 0, 0);
    const pLensHit = new THREE.Vector3(0, objectHeight, 0);
    const pFocusBack = new THREE.Vector3(0, 0, focalLength);
    const pImgTip = new THREE.Vector3(0, imageHeight, imageDist);

    // รังสีที่ 1: ขนานแกนมุขสำคัญ แล้วหักเหผ่านจุดโฟกัสด้านหลังสู่ปลายภาพ
    const ray1Points = [pObjTip, pLensHit, pImgTip];
    // รังสีที่ 2: วิ่งตรงผ่านกึ่งกลางเลนส์โดยไม่หักเห
    const ray2Points = [pObjTip, pLensCenter, pImgTip];

    return {
        imageDist: imageDist.toFixed(2),
        magnification: magnification.toFixed(2),
        isReal: imageDist > 0,
        ray1: ray1Points,
        ray2: ray2Points
    };
}
\end{codebox}

\section{ระบบส่งออกรายงานผลการทดลองวิชาการอัตโนมัติ}

ในกระบวนการเรียนรู้เชิงรุก (Active Learning) สิ่งสำคัญไม่ยิ่งหย่อนไปกว่าการทดลองคือการบันทึกและสังเคราะห์ผลการทดลองเป็นรายงานวิชาการ เพื่ออำนวยความสะดวกแก่นักเรียน สถาปัตยกรรมห้องปฏิบัติการเสมือนจริงได้ผสาน \textbf{ระบบส่งออกรายงานการทดลองอัตโนมัติ (Automated Lab Report Document Export)}

ระบบจะเก็บรวบรวมข้อมูลตัวแปรต้น ตัวแปรตาม และค่าสถิติความคลาดเคลื่อนตลอดการทดลองของผู้เรียน แล้วสร้างเอกสารสรุปผลการทดลองในรูปแบบ HTML หรือ Markdown ที่พร้อมแปลงเป็น PDF ทันทีบนฝั่งไคลเอนต์ (Client-Side Document Assembly) ดังโครงสร้างข้อมูล:
\begin{enumerate}
    \item \textbf{ข้อมูลประจำตัวและการทดลอง:} ชื่อผู้เรียน รหัสนักศึกษา วันที่ และชื่อโมดูลการทดลอง
    \item \textbf{ตารางผลการวัดเชิงปริมาณ (Data Tables):} ค่าตัวแปรที่บันทึก เช่น ระยะวัตถุ $s$, ระยะภาพ $s'$, แรงเสียดทาน, อุณหภูมิ
    \item \textbf{การวิเคราะห์ความคลาดเคลื่อน (Error Analysis):} ค่าเฉลี่ย ส่วนเบี่ยงเบนมาตรฐาน และเปอร์เซ็นต์ความคลาดเคลื่อน
    \item \textbf{ภาพถ่ายบันทึกหลักฐาน (In-App Snapshot):} ภาพสแน็ปช็อตสามมิติขณะทำการทดลองจากมุมมองของผู้เรียน
\end{enumerate}

\section{กิจกรรมเชิงรุก AR STEM Scavenger Hunt}

นอกเหนือจากการทดลองภายในห้องแล็บเสมือนแบบปิด การจัดกิจกรรมการเรียนรู้แบบ \textbf{AR STEM Scavenger Hunt} ช่วยเชื่อมโยงผู้เรียนเข้ากับโลกกายภาพจริง โดยใช้สมาร์ตโฟนหรือแท็บเล็ตร่วมกับเทคโนโลยี Augmented Reality (AR)

ผู้เรียนจะต้องออกสำรวจพื้นที่จริงในมหาวิทยาลัยหรือโรงเรียนเพื่อค้นหา "มาร์กเกอร์วิทยาศาสตร์" ที่ซ่อนอยู่ เมื่อสแกนพบ ระบบจะเรนเดอร์คำถามสะเต็ม โจทย์ฟิสิกส์ หรือการทดลอง AR ทับซ้อนลงบนวัตถุจริง เช่น การวัดความเร่งของลิฟต์โดยสาร หรือการวิเคราะห์การสะท้อนของแสงบนกระจกอาคาร กิจกรรมนี้ส่งเสริมทักษะการทำงานเป็นทีม การแก้ปัญหาเฉพาะหน้า และการคิดวิเคราะห์ในสถานการณ์จริง

\clearpage

\section*{สรุปท้ายบทที่ 7}
\addcontentsline{toc}{section}{สรุปท้ายบทที่ 7}

การพัฒนาห้องปฏิบัติการเสมือนจริงขั้นสูงด้วยมาตรฐาน OpenXR และการประมวลผลเชิงพื้นที่เป็นการหลอมรวมองค์ความรู้ด้านวิทยาศาสตร์ คณิตศาสตร์ การเขียนโปรแกรม และการออกแบบปฏิสัมพันธ์เข้าด้วยกันอย่างสมบูรณ์ ตั้งแต่การจำลองทัศนศาสตร์ การหมุนเฟสเซอร์ในวงจรไฟฟ้ากระแสสลับ ไปจนถึงการจัดกิจกรรม AR STEM Scavenger Hunt และระบบส่งออกรายงานอัตโนมัติ ช่วยทลายข้อจำกัดของการทดลองในห้องแล็บจริง เปิดโอกาสให้นักเรียนทุกคนสามารถเข้าถึงการทดลองทางวิทยาศาสตร์ระดับสูงได้อย่างเท่าเทียม ปลอดภัย และเกิดการเรียนรู้อย่างยั่งยืน

\clearpage

\section*{แบบฝึกหัดท้ายบทที่ 7}
\addcontentsline{toc}{section}{แบบฝึกหัดท้ายบทที่ 7}

\begin{enumerate}
    \item จงออกแบบขั้นตอนวิธีในการคำนวณตำแหน่งและทิศทางของรังสีแสงหักเหผ่านเลนส์เว้าในปริภูมิสามมิติ
    \item อธิบายความหมายทางกายภาพของจุดลากรองจ์ $L_1$ ในระบบดาวคู่ และยกตัวอย่างว่าการจำลองสามมิติช่วยให้เข้าใจปรากฏการณ์นี้ดีขึ้นอย่างไร
    \item เสนอแนวทางการออกแบบห้องปฏิบัติการเสมือนจริงสำหรับโรงเรียนที่ขาดแคลนอุปกรณ์ฮาร์ดแวร์แว่น XR ราคาแพง
    \item ในวงจรอนุกรม $RLC$ กำหนดให้ $R = 30\,\Omega$, $L = 0.1\text{ H}$, และ $C = 50\,\mu\text{F}$ ต่อกับแหล่งจ่ายไฟ $v(t) = 100\cos(1000 t)\text{ V}$ จงคำนวณอิมพีแดนซ์รวม $Z$, มุมต่างเฟส $\phi$, และความถี่เรโซแนนซ์ $f_0$
    \item อธิบายหลักการทำงานของระบบส่งออกรายงานผลการทดลองอัตโนมัติบนฝั่งไคลเอนต์ (Client-Side Lab Report Export) และประโยชน์ที่มีต่อการประเมินผลการเรียนรู้เชิงรุก
\end{enumerate}
